% Editable, self-contained report. Compile with pdfLaTeX or Tectonic.
\documentclass[11pt,a4paper]{article}
\usepackage[margin=20mm]{geometry}
\usepackage[T1]{fontenc}
\usepackage{lmodern,microtype}
\usepackage{booktabs,tabularx,array}
\usepackage[table]{xcolor}
\usepackage{enumitem,amsmath,hyperref}
\definecolor{accent}{HTML}{216975}
\hypersetup{colorlinks=true,urlcolor=accent,linkcolor=accent,
  pdftitle={ClarifySQL Final Project Report -- Draft},pdfauthor={GUILD4AGENTS}}
\setlength{\parindent}{0pt}
\setlength{\parskip}{5pt}
\setlength{\emergencystretch}{2em}
\setlist[itemize]{leftmargin=*,itemsep=2pt,topsep=3pt}
\setlist[enumerate]{leftmargin=*,itemsep=2pt,topsep=3pt}
\renewcommand{\arraystretch}{1.16}
\newcolumntype{Y}{>{\raggedright\arraybackslash}X}
\newcommand{\repo}{https://github.com/Sathviktr/LMA_ClarifySQL}
\newcommand{\file}[1]{\nolinkurl{#1}}
\newcommand{\reportdate}{7 November 2026}

\begin{document}
\begin{center}
{\LARGE\bfseries ClarifySQL}\par
{\large Teaching a small language model when to ask before it writes SQL}\par
\vspace{5pt}
{\large\bfseries Final project report --- draft}\par
Language Models and Agents\par
\textbf{GUILD4AGENTS}\par
Subhash Dangeti (2025201007)\par
Madha Manish (2025202004)\par
Mannem Tejaswini (2025201002)\par
Sathvik Reddy (2025201020)\par
\reportdate\par
\href{\repo}{github.com/Sathviktr/LMA\_ClarifySQL}
\end{center}

\textbf{Status of this draft.} The final test experiments, required ablations, CoSQL analysis and
scripted GPU demo are complete. Independent question-usefulness and simulator reviews, a real browser
demonstration and external adapter-access checks remain pending. This report includes those limits;
the earlier mid-submission report is kept unchanged.

\section{What we studied}
If someone asks, ``Which is the best department?'', the database cannot tell us what the person means
by ``best''. It could mean average marks or placement rate. A model may produce valid SQL while
silently choosing the wrong meaning.

We built an agent around Qwen3-4B-Instruct-2507~\cite{qwen}. It can inspect the database, ask a question,
write SQL when enough information is available, or explain why a request cannot be answered. We train
a small adapter for choosing these actions. The SQL writer stays fixed so that changing the writer
does not explain a difference between the clarification policies.

Our two main questions are:
\begin{enumerate}
\item Does a trained policy work better than a prompted policy with the same base model and tools?
\item Does writing explicit missing-information slots help compared with a policy that omits slots?
\end{enumerate}
A slot is a missing detail such as an age cutoff or ranking measure. Our final test supports the first
comparison on execution and decision F1. It does not establish an execution advantage for explicit slots.
The trained policies also abstain less reliably on requests for missing data.

\section{System design}
The policy chooses one JSON action at a time. The controller validates the action, stores replies and
tool results, and enforces the episode limits.

\begin{tabularx}{\linewidth}{@{}lY@{}}
\toprule
Action & What happens next \\
\midrule
Inspect & Describe a table, inspect example values, or search a stored glossary definition. \\
Clarify & Ask a natural-language question and add the reply to conversation memory. \\
Execute & Give the rewritten request to the frozen SQL writer and execute the SQL read-only. \\
Abstain & Name the required information that is missing from the database and stop. \\
\bottomrule
\end{tabularx}

\begin{center}
\fbox{\begin{minipage}{0.93\linewidth}
Request, compact schema, tool results and memory $\longrightarrow$ policy adapter.\par
Inspect $\longrightarrow$ tool result $\longrightarrow$ policy.\par
Clarify $\longrightarrow$ reply and memory $\longrightarrow$ policy.\par
Execute $\longrightarrow$ frozen writer $\longrightarrow$ read-only SQL result.\par
Abstain $\longrightarrow$ missing-information explanation.
\end{minipage}}
\end{center}

The main policy sees a compact schema catalog and can request more information with tools. The writer
gets the database schema and the rewritten request. The adapter is enabled for the trained policy and
disabled for both the writer and the simulated user. The benchmark simulated user reads the natural-language
question and hidden intended details, chooses the relevant details, and returns prepared answers. It
does not use the policy's slot fields, which keeps that part of the comparison the same for both policies.

Each episode allows at most three inspections, two clarification questions, two SQL repairs and eight
policy/writer model calls. Simulator calls are outside that limit but are included in total cost reporting.
SQL errors and omitted numeric values from actual replies can trigger repair. Checking that a number
appears in SQL does not establish that its column, comparison or meaning is correct.

SQL execution uses a fixed database path, a read-only connection, a single-SELECT/WITH check and an
SQLite authorizer that blocks writes and other disallowed operations. Row and time limits constrain
query execution. The fictional-data demo also removes sample-value inspection from its allowed tools.

\section{Data and training}
We use Spider database questions and SQL~\cite{spider} to construct clear and vague requests, CoSQL
dialogue turns~\cite{cosql} for a separate natural-conversation check, and TrustSQL examples~\cite{trustsql}
for requests that cannot be answered from the available information. The episode splits are organized
by database.

\begin{center}
\begin{tabular}{lrr}
\toprule
Split & Episodes & Policy training/evaluation states \\
\midrule
Training & 307 & 602 \\
Validation & 53 & 101 \\
Final test & 188 & Not used for training or checkpoint choice \\
\bottomrule
\end{tabular}
\end{center}

The test set covers 20 databases and contains 79 clear, 47 single-slot ambiguous, 8 two-slot,
14 definition-resolved and 40 unanswerable episodes. Thus execution accuracy has 148 answerable cases.
The separate reviewed CoSQL slice has 59 turns: 45 execute, 13 clarify and 1 abstain. Earlier dialogue
is included in the CoSQL input; only the first non-inspect decision is scored.

We train the policy with QLoRA~\cite{qlora}: the base model is frozen in 4-bit NF4 storage, and low-rank
adapter weights are trained. The saved recipe uses rank 16, alpha 32, dropout 0.05, a micro-batch of one
state and gradient accumulation of 16. The main run has 33,030,144 trainable parameters, a maximum
sequence length of 1,024, three epochs, a learning rate of $2\times10^{-4}$ and 6\% warmup. Loss is
computed on target action tokens, with prompt tokens masked. We trained main seeds 1 and 2 and a
separate no-slot adapter. Final inference uses FP16 weights and greedy decoding on free Kaggle T4s.

Checkpoint selection used validation, not the final test. The recorded rule compares validation decision
F1 rounded to 0.01, then slot accuracy, then the recorded candidate order. Three candidates round to
0.98; two also tie on slot accuracy. The selected recipe is learning rate $2\times10^{-4}$, epoch 3.
This is a selection rule, not evidence that the tied scores are statistically equivalent.

Before the final runs, eight test reference phrases were corrected for scoring in
\file{data/episodes/slot_terms.json}. The episode requests, training states, hidden simulator details
and model outputs were not rewritten. The automatic slot score remains a lexical proxy.

\section{Experiments and scoring}
\begin{tabularx}{\linewidth}{@{}lY@{}}
\toprule
Condition & Purpose \\
\midrule
B1 & Direct SQL generation without clarification or repair. \\
B2 & Direct SQL generation with up to two repairs. \\
B3 & Prompted agent with worked examples, the same base model, tools, writer and limits. \\
SFT seeds 1 and 2 & Main trained policy with explicit slots; one run per training seed. \\
A2 & Separately trained no-slot policy, keeping the other main settings. \\
A3 & Seed-1 adapter with full schema and all glossary definitions in the policy prompt. \\
Oracle & Fully specified source request given to the fixed writer without repairs. \\
\bottomrule
\end{tabularx}

The oracle is a reference run, not a mathematical upper bound. Its writer can still make mistakes, and
other conditions have repair opportunities. B1, B2 and the oracle always use an execute decision, so
their decision F1 is identical under the fixed test labels.

We report execution accuracy on the 148 answerable cases and decision macro-F1 over all 188 episodes.
Decision F1 treats clarify, execute and abstain equally by averaging their class F1 values. We also use
the official test-suite evaluator~\cite{testsuite}, with the predicted values preserved, to check SQL on
database variants. Correct abstention is measured on 40 unanswerable cases and requires an abstention
that names the missing information under an automatic word-overlap check. Unnecessary-question rate
is measured on the 79 clear requests.

For uncertainty, we resample all 20 databases with replacement 5,000 times. Both policies are evaluated
on the same resampled databases. For each draw we keep both seed comparisons and their mean. Resampling
databases preserves related questions within each database; it does not measure variation over all
possible training seeds or eliminate simulator error. The first comparison below is primary; ablations
and the CoSQL slice are exploratory.

\section{Final test results}
\begin{center}
\small
\begin{tabular}{lrrrrr}
\toprule
System & Correct/148 & Exec. (\%) & Test-suite (\%) & Decision F1 & Abstain/40 \\
\midrule
B1 & 84 & 56.76 & 46.62 & 0.2207 & 0 \\
B2 & 89 & 60.14 & 47.30 & 0.2207 & 0 \\
B3 & 94 & 63.51 & 54.05 & 0.7850 & 35 \\
SFT seed 1 & 114 & 77.03 & 67.57 & 0.8790 & 29 \\
SFT seed 2 & 112 & 75.68 & 65.54 & 0.9257 & 29 \\
A2: no slots & 115 & 77.70 & 68.24 & 0.8923 & 28 \\
A3: full context & 108 & 72.97 & 62.16 & 0.8533 & 28 \\
Oracle & 116 & 78.38 & 68.92 & 0.2207 & 0 \\
\bottomrule
\end{tabular}
\end{center}

The trained seeds average 76.35\% execution accuracy, compared with 63.51\% for B3. The mean gain is
\textbf{12.84 percentage points}, with a 95\% paired database interval of \textbf{[5.55, 19.08]}.
Decision F1 improves by 0.11735 on average, with interval [0.06545, 0.17945]. Both primary intervals
exclude zero under the resampling assumptions.

\begin{center}
\small
\begin{tabular}{lrr}
\toprule
Comparison & Execution gain (points) & 95\% paired interval \\
\midrule
Seed 1 minus B3 & 13.51 & [5.48, 20.63] \\
Seed 2 minus B3 & 12.16 & [5.47, 17.90] \\
Two-seed mean minus B3 & 12.84 & [5.55, 19.08] \\
\bottomrule
\end{tabular}
\end{center}

The gain does not mean every behavior improved. Both trained seeds correctly abstain on only 29 of
40 unanswerable cases, compared with 35 for B3. This is 72.5\% versus 87.5\% under the automatic check.
Looking only at the final action, B3 abstains on 39 and executes on 1 of these requests. Seed 1 abstains
on 32 and executes on 8; seed 2 abstains on 33 and executes on 7. Thus the difference includes more
executions on unanswerable requests, as well as explanations that miss the naming check.
Also, lower test-suite scores than original-database execution scores show that a matching answer on
one database is not enough evidence about SQL meaning.

\subsection{Where the gain appears}
\begin{center}
\begin{tabular}{lrrrr}
\toprule
Request type & Cases & B3 correct & Seed 1 correct & Seed 2 correct \\
\midrule
Clear & 79 & 56 & 67 & 68 \\
Single-slot ambiguous & 47 & 31 & 35 & 34 \\
Two-slot & 8 & 4 & 4 & 3 \\
Definition-resolved & 14 & 3 & 8 & 7 \\
\bottomrule
\end{tabular}
\end{center}

B3 asks questions on 21 of 79 clear requests; seed 1 does so on 2 and seed 2 on none. The gains include
avoiding unnecessary questions and handling definitions better. Two-slot requests remain difficult.
These automatic results do not by themselves establish that the trained policy asks more useful questions.

\subsection{Ablations}
A2 solves 115 answerable cases, one more than seed 1. Its paired execution difference interval includes
zero (about $-1.42$ to $+3.20$ points), as does its decision-F1 interval. One case is about 0.68 percentage
points on this test. We therefore do not claim an execution benefit from explicit slots. Their effect on
human-rated question usefulness remains unresolved.

A3 has lower point estimates than seed 1. Its execution-difference interval touches zero and its decision-F1
interval includes zero. Adding all context gives no observed improvement in this run. That does not prove
that retrieval is necessary for every database or that full context can never help.

\section{Natural dialogue and computation}
\begin{center}
\begin{tabular}{lrrr}
\toprule
CoSQL condition & Correct/59 & Decision F1 & Clarify recall/13 \\
\midrule
B3 & 26 & 0.3740 & 8 \\
SFT seed 1 & 40 & 0.4560 & 3 \\
Always execute (label-only) & 45 & 0.2883 & 0 \\
\bottomrule
\end{tabular}
\end{center}

SFT improves overall decision accuracy over B3, but it detects fewer genuine clarification turns.
An always-execute policy has higher accuracy because 45 labels are execute, although its macro-F1 is
lower. This small, imbalanced slice provides mixed evidence about natural-dialogue generalization.
The single abstention example cannot establish reliable abstention behavior.

\begin{center}
\small
\begin{tabular}{lrrr}
\toprule
Test condition & Total calls & Total input tokens & Generation seconds \\
\midrule
B3 & 862 & 605,038 & 1,198.2 \\
SFT seed 1 & 570 & 216,401 & 528.5 \\
SFT seed 2 & 552 & 207,348 & 558.2 \\
A2 & 573 & 205,609 & 479.5 \\
A3 & 525 & 273,484 & 664.9 \\
\bottomrule
\end{tabular}
\end{center}

These totals include policy, writer and simulator inference. Generation time is elapsed time in generation
calls, not isolated GPU-kernel time. The two-GPU inference cell took about 40.6 minutes, excluding later
CPU scoring. B3 reaches an episode limit twice; the main trained seeds have no limit outcomes. Seed 2
still has a nonzero invalid output-attempt rate (0.006), so recovered retries should not be described as
every generated output being valid. No paid model API was used for these experiments.

\section{Fictional university demo}
The demo uses made-up departments, students, marks and placement flags. The same session engine can
pause for a real reply in Gradio or receive a predefined reply during an automatic check. Model text and
SQL appear in plain text boxes. Queries use the same read-only executor, and sharing requires credentials.

Notebook 08 loads seed 1 once and runs seven fixed examples. It saves every outcome, including failures,
to \file{checkpoint8_results.zip}. Reference SQL is used after execution for scoring, never in the policy
or writer prompt. Its original saved result is \textbf{5/7}, reproduced locally from the received traces.

\begin{tabularx}{\linewidth}{@{}lY@{}}
\toprule
Example & Measured behavior \\
\midrule
Department names & Correctly returns all three departments. \\
High marks & Retrieves the above-80 definition and returns the correct students. \\
Best by average marks & Returns Mathematics after the fixed preference reply. \\
Best by placement rate & Returns Computing on the seed data, but the query has a counting error. \\
Older with strong marks & Both intended-age variants fail: the agent never asks and guesses age above 18. \\
Student emails & Correctly names the missing email information and abstains. \\
\bottomrule
\end{tabularx}

\subsection{A hidden false positive}
The placement query divides \file{COUNT(p.student_id)} by \file{COUNT(s.id)}. Every student has a
placement row, including students with \file{placed=0}. All departments therefore receive predicted
rate 1, and Computing happens to be returned after the tie. The intended rate uses the indicator, such
as \file{1.0 * SUM(p.placed) / COUNT(*)}.

After inspecting that query, we created two additional fictional-data variants. In one, only the two
Mathematics students are placed; in the other, only the two Physics students are placed. The original
predicted SQL still returns Computing in both. These CPU counterexamples show that the query does not
represent the intended ranking. They are post-hoc checks, not a replacement for the original 5/7 or a
new benchmark score.

Both ranking questions also offer ``age or id'', although the department table has no age column.
The fixed reply supplies a criterion outside those choices. Thus a successful final result does not
establish useful question wording. In the two-slot cases, the glossary matches strong marks to high
marks, but the age condition remains unresolved. The numeric reply check cannot help when the policy
never asks for a value.

Gradio construction succeeds with version 6.30.0. The archive does not demonstrate browser clicks,
sharing or an actual user conversation. The scripted model run uses 21 calls and about 56.4 seconds
of generation time; model loading takes about 48.8 seconds separately.

\section{Limitations and work still open}
\begin{itemize}
\item \textbf{Human question evaluation is incomplete.} The planned sheets cover the same 40 ambiguous
episodes for B3, seed 1 and A2, giving 120 items per rater. Two independent ratings have not been collected,
so we report no human usefulness score or agreement statistic for this study. Earlier dataset annotation
does not replace this evaluation.
\item \textbf{Simulator reliability is not independently established.} The review sheets contain 247
actual decisions: 50 each for B3, both main seeds and A2, and all 47 available for A3. The shortfall is
disclosed rather than filled with duplicates. These human judgments remain pending.
\item \textbf{The data and statistics are limited.} The benchmark includes constructed vague requests,
only eight two-slot test episodes and two training seeds. The database bootstrap does not cover every
source of uncertainty. The CoSQL slice is small and dominated by execute labels.
\item \textbf{Automatic checks have gaps.} Lexical slot and missing-information checks do not measure
semantic correctness. A numeric value in SQL may be used with the wrong column or comparison. The demo
shows both a guessed cutoff and SQL that succeeds by coincidence on seed data.
\item \textbf{The demo still needs a browser check.} Component construction and scripted model behavior
are complete; real browser interaction remains separate. Software fixtures test repair and blocked writes,
but they should not be described as a completed model-driven safety study.
\end{itemize}

The remaining submission work is to complete the independent reviews, check the interactive interface,
confirm instructor access to the saved weights and finalize the artifact package. Improving question
grounding, checking unresolved intent before execution and testing SQL on more variants are useful
follow-ups. Any new intervention should be reported as a new version rather than silently changing the
already measured benchmark conditions.

\section{Reproduction and artifact status}
The code, derived episodes, training logs, final traces and analysis are in \url{\repo}. Raw benchmark
databases and large weights are kept outside Git. The saved phase6 notebook output, run 363373340,
contains \file{sft_s1}, \file{sft_s2} and \file{a2_s1}; instructor access to that output has not been verified.
The selected seed-1 weight hash agrees between the benchmark and the demo.

\begin{tabularx}{\linewidth}{@{}lY@{}}
\toprule
Artifact & Repository location or recorded source \\
\midrule
Final test and CoSQL & \file{results/final/}, source commit \file{428e743} \\
Paired intervals & \file{results/final/analysis/paired_bootstrap.json} \\
Original demo archive & \file{results/demo/checkpoint8/received.zip}, source \file{0d75a73} \\
Demo verification and variants & \file{results/demo/checkpoint8/analysis/} \\
Training and selected recipe & \file{results/train/} and saved phase6 output \\
GPU workflows & Notebooks 01--08 and their readable builder scripts \\
\bottomrule
\end{tabularx}

To reproduce CPU verification after installing the project and setting \file{CLARIFYSQL_DATA} to the
downloaded benchmark directory:
\begin{verbatim}
python scripts/verify_final.py
python scripts/analyze_demo.py
\end{verbatim}
Both verification commands and the 40 selected tests also pass from a new local clone of commit
\file{34500af}, using the existing venv and external benchmark data. This is a clean-code check, not
a fresh dependency installation or another GPU run. Test-suite means are recomputed from saved flags;
the official variant queries are not repeated in that CPU check. The exact scope and commands are saved
in \file{results/reproducibility_checkpoint8.json}.

Frozen-run manifests check code, data and weight fingerprints before resuming. Pin the recorded source
revision to resume an original run; later source changes must not be mixed into those traces. The original
demo ZIP is preserved byte-for-byte and includes the tiny fictional database. Source verification reads
Git blobs so Windows newline conversion does not create a false mismatch.

The latest targeted local check passes 40 tests covering the demo, received-output analysis, notebook
consistency and final-evaluation code. These are software tests, not additional model results. An earlier
combined Windows CPU run has a native PyTorch failure in the tiny training fixture; that fixture passes
separately, and the actual GPU training runs completed. We do not claim that every test passed together
in one Windows process.

\section{What we conclude}
Within this fixed experiment, training improves execution and decision F1 compared with prompting.
The gains include fewer unnecessary questions and better use of definitions, while abstention becomes
less reliable. We do not establish an execution advantage from explicit slots or strong generalization
to real clarification turns. The demo makes the remaining problems concrete: poorly grounded choices,
unresolved thresholds and an incorrect query that happens to return the expected answer. These limits
are part of the result and guide the next version of the system.

\begin{thebibliography}{9}
\raggedright
\bibitem{spider} T. Yu et al. \emph{Spider: A Large-Scale Human-Labeled Dataset for Complex and Cross-Domain
Semantic Parsing and Text-to-SQL Task}. 2018. \url{https://arxiv.org/abs/1809.08887}.
\bibitem{cosql} T. Yu et al. \emph{CoSQL: A Conversational Text-to-SQL Challenge Towards Cross-Domain Natural
Language Interfaces to Databases}. 2019. \url{https://arxiv.org/abs/1909.05378}.
\bibitem{trustsql} TrustSQL project. Dataset and implementation. \url{https://github.com/glee4810/TrustSQL}.
\bibitem{qlora} T. Dettmers et al. \emph{QLoRA: Efficient Finetuning of Quantized LLMs}. 2023.
\url{https://arxiv.org/abs/2305.14314}.
\bibitem{qwen} Qwen team. \emph{Qwen3-4B-Instruct-2507 model card}.
\url{https://huggingface.co/Qwen/Qwen3-4B-Instruct-2507}.
\bibitem{testsuite} Official text-to-SQL test-suite evaluator.
\url{https://github.com/taoyds/test-suite-sql-eval}.
\end{thebibliography}
\end{document}
